\documentclass[../book.tex]{subfiles}
\begin{document}

    \section{数列极限}

    \subsection{数列的概念}

    按一定次序排列的一列数$a_1, a_2, \cdots, a_n, \cdots$称为数列，简记为$\{a_n\}$。其中$a_n$称为数列的\textbf{通项}（或一般项）。若数列$\{a_n\}$的第$n$项$a_n$与$n$之间可以用一个公式表示，则该公式称为数列的\textbf{通项公式}。

    \begin{enumerate}
        \item \textbf{子列}：从数列$\{a_n\}$中任意选取无穷多项，保持原有顺序组成的新数列，记为$\{a_{n_k}\}$，其中$n_1 < n_2 < \cdots < n_k < \cdots$

        \item \textbf{等差数列}：$a_n = a_1 + (n-1)d$，其中$a_1$为首项，$d$为公差
        \begin{itemize}
            \item 通项公式：$a_n = a_1 + (n-1)d$
            \item 前$n$项和：$S_n = \dfrac{n(a_1 + a_n)}{2} = \dfrac{n[2a_1 + (n-1)d]}{2}$
        \end{itemize}

        \item \textbf{等比数列}：$a_n = a_1 q^{n-1}$，其中$a_1$为首项，$q$为公比（$q \neq 0$）
        \begin{itemize}
            \item 通项公式：$a_n = a_1 q^{n-1}$
            \item 前$n$项和：$S_n = \begin{cases}
                                       na_1, & q = 1 \\ a_1 \cdot \dfrac{1-q^n}{1-q}, & q \neq 1
            \end{cases}$
            \item $1 + r + r^2 + \cdots + r^{n-1} = \dfrac{1-r^n}{1-r}(r \neq 1)$
        \end{itemize}

        \item \textbf{单调数列}：
        \begin{itemize}
            \item 单调不减：$a_{n+1} \ge a_n$（单调递增：$a_{n+1} > a_n$）
            \item 单调不增：$a_{n+1} \le a_n$（单调递减：$a_{n+1} < a_n$）
        \end{itemize}

        \item \textbf{有界数列}：若存在$M > 0$，使得对所有$n$都有$|a_n| \le M$，则称$\{a_n\}$有界
        \begin{itemize}
            \item 有上界：存在$M$，使得$a_n \le M$
            \item 有下界：存在$m$，使得$a_n \ge m$
            \item 证明数列有界的几种方法
            \begin{circlenum}
                \item 找$M$，使得$|a_n| \leq M$
                \item 放缩法
                \item 找最值
                \item 基本不等式法
            \end{circlenum}
        \end{itemize}

        \item \textbf{一些常见数列前$n$项的和}：
        \begin{itemize}
            \item $\displaystyle\sum_{k=1}^{n} k = 1 + 2 + 3 + \cdots + n = \dfrac{n(n+1)}{2}$
            \item $\displaystyle\sum_{k=1}^{n} k^2 = 1^2 + 2^2 + 3^2 + \cdots + n^2 = \dfrac{n(n+1)(2n+1)}{6}$
            \item $\displaystyle\sum_{k=1}^{n} \dfrac{1}{k(k+1)} = \dfrac{1}{1 \cdot 2} + \dfrac{1}{2 \cdot 3} + \cdots + \dfrac{1}{n(n+1)} = \dfrac{n}{n+1}$（裂项相消）
        \end{itemize}

        \item \textbf{一个重要数列$\{(1 + \dfrac{1}{n})^n\}$的结论}：
        \begin{itemize}
            \item 单调递增且有上界
            \item $\lim\limits_{n \to \infty} \left(1 + \dfrac{1}{n}\right)^n = e$
            \item 更一般地：$\lim\limits_{n \to \infty} \left(1 + \dfrac{a}{n}\right)^n = e^a$
        \end{itemize}
    \end{enumerate}

    \subsection{数列极限的定义}

    设$\{x_n\}$为一数列，若存在常数$a$，对于任意的$\epsilon > 0$（不论它多么小），总存在正整数$N$，使得当$n > N$时，$|x_n - a| < \epsilon$，则称常数$a$是数列$\{x_n\}$的\textbf{极限}，或者称数列$\{x_n\}$\textbf{收敛于}$a$，记为
    \begin{gather*}
        \lim_{n \to \infty} x_n = a \quad \text{或} \quad x_n \to a \;(n \to \infty)
    \end{gather*}

    \textbf{$\epsilon-N$语言}：
    \[
        \forall \epsilon > 0,\ \exists N > 0,\ n > N \Rightarrow |x_n - a| < \epsilon
    \]

    \begin{enumerate}
        \item \textbf{常用的语言}
        \begin{itemize}
            \item $\lim\limits_{n\to \infty} x_n = a \Leftrightarrow \forall \epsilon > 0,\ \exists N > 0,\ n > N \Rightarrow |x_n - a| < \epsilon$
            \begin{circlenum}
                \item 与函数极限的定义做对比：$\lim\limits_{x\to \infty} f(x) = A \Leftrightarrow \forall \epsilon > 0,\ \exists X > 0,\ |x| > X \Rightarrow |f(x) - A| < \epsilon$
                \item 定义中的 $n > N$，$N$未必只能取整数
            \end{circlenum}
            \item $\lim\limits_{n\to \infty} x_n = \infty \Leftrightarrow \forall M > 0,\ \exists N > 0,\ n > N \Rightarrow |x_n| > M$
            \item $\lim\limits_{n\to \infty} x_n = +\infty \Leftrightarrow \forall M > 0,\ \exists N > 0,\ n > N \Rightarrow x_n > M$
            \item $\lim\limits_{n\to \infty} x_n = -\infty \Leftrightarrow \forall M > 0,\ \exists N > 0,\ n > N \Rightarrow x_n < -M$
        \end{itemize}
        \item \textbf{几何解释}：将常数$a$及数列$x_1, x_2, \cdots, x_n, \cdots$在数轴上表示出来。"$\lim\limits_{n\to\infty}x_n = a$"的几何意义是：对于任意$\epsilon > 0$，存在正整数$N$，使得当$n > N$时，所有点$x_n$都落在$(a-\epsilon, a+\epsilon)$内，而在这个邻域外至多有$N$个点。
    \end{enumerate}

    \subsection{收敛数列的性质}
    \begin{enumerate}
        \item \textbf{唯一性}：给出数列$\{x_n\}$，若$\lim\limits_{n\to\infty} x_n = a$（存在），则$a$是唯一的

        \item \textbf{有界性}：若数列$\{x_n\}$收敛，则$\{x_n\}$有界
        \begin{itemize}
            \item \textbf{注}：有界是收敛的必要条件，但非充分条件
            \item 反例：$x_n = (-1)^n$有界但不收敛
        \end{itemize}

        \item \textbf{保号性}：设$\lim\limits_{n\to\infty} x_n = a$
        \begin{enumerate}
            \item 若$a > b$，则存在正整数$N$，使得当$n > N$时，$x_n > b$
            \item 若数列 $\{x_n\}$ 从某项起有 $x_n \geq b$，则 $a \geq b$，其中 $b$ 为任意实数，常考 $b = 0$ 的情形
        \end{enumerate}

        \item \textbf{脱帽与带帽解法}（常考情形为$b = 0$或与$0$比较）：
        \begin{itemize}
            \item \textbf{脱帽法（严格不等号）}：极限符号去掉后不等号方向不变
            \[
                \lim_{n\to\infty} x_n > b \quad \Rightarrow \quad x_n > b \quad (n > N)
            \]
            \[
                \lim_{n\to\infty} x_n < b \quad \Rightarrow \quad x_n < b \quad (n > N)
            \]

            \item \textbf{带帽法（非严格不等号）}：加上极限符号后不等号可能变弱
            \[
                x_n \ge b \;(n > N) \quad \Rightarrow \quad \lim_{n\to\infty} x_n \ge b
            \]
            \[
                x_n \le b \;(n > N) \quad \Rightarrow \quad \lim_{n\to\infty} x_n \le b
            \]
            \item \textbf{注}：$x_n > b$只能推出$\lim\limits_{n\to\infty} x_n \ge b$，严格不等号可能变弱。\emph{e.g.} $x_n = 1 + \dfrac{1}{n} > 1$，但$\lim\limits_{n\to\infty} x_n = 1$
        \end{itemize}
        \item \textbf{数列收敛与其子列收敛的关系}：
        \begin{itemize}
            \item 若数列$\{x_n\}$收敛于$a$，则其任意子列$\{x_{n_k}\}$也收敛于$a$
            \item 若数列$\{x_n\}$的某个子列发散，则$\{x_n\}$发散
            \item 若数列$\{x_n\}$有两个子列收敛于不同的极限，则$\{x_n\}$发散
        \end{itemize}
    \end{enumerate}

    \subsection{极限四则运算规则}

    设$\lim\limits_{n\to\infty} x_n = a$，$\lim\limits_{n\to\infty} y_n = b$，则：

    \begin{enumerate}
        \item 使用四则运算前，必须确保各数列极限存在
        \item $\lim\limits_{n\to\infty} (x_n \pm y_n) = a \pm b$

        \item \textbf{乘法}：$\lim\limits_{n\to\infty} (x_n \cdot y_n) = a \cdot b$

        \item \textbf{数乘}：$\lim\limits_{n\to\infty} (c \cdot x_n) = c \cdot a$（$c$为常数）

        \item \textbf{除法}：若$y_n \neq 0$且$b \neq 0$，则$\lim\limits_{n\to\infty} \dfrac{x_n}{y_n} = \dfrac{a}{b}$

        \item \textbf{幂运算}：$\lim\limits_{n\to\infty} x_n^k = a^k$（$k$为正整数）
    \end{enumerate}

    \subsection{海涅定理（归结原则）}

    \begin{enumerate}
        \item \textbf{定理内容}：设$f(x)$在$\mathring{U}(x_0, \delta)$内有定义，则$\lim\limits_{x\to x_0} f(x) = A$存在 $\Leftrightarrow$ 对任何$\mathring{U}(x_0, \delta)$内以$x_0$为极限的数列$\{x_n\}(x_n \neq x_0)$，极限$\lim\limits_{n\to\infty} f(x_n) = A$存在

        \item \textbf{核心思想}：在极限存在的条件下，\textbf{函数极限和数列极限可以相互转化}

        \item \textbf{常用情形}：
        \begin{itemize}
            \item 当$x\to 0$时，取$x_n = \dfrac{1}{n}$，若$\lim\limits_{x\to 0} f(x) = A$，则$\lim\limits_{n\to\infty} f\left(\dfrac{1}{n}\right) = A$
            \item 当$x\to +\infty$时，取$x_n = n$，若$\lim\limits_{x\to +\infty} f(x) = A$，则$\lim\limits_{n\to\infty} f(n) = A$
            \begin{analysisbox}[Example]
                \[
                    \lim_{n\to\infty}\frac{\ln n}{n}
                    = \lim_{x\to+\infty}\frac{\ln x}{x}
                    \overset{\text{洛必达}}=
                    \lim_{x\to+\infty}\frac{\displaystyle\frac{1}{x}}{1}
                    = 0
                \]
            \end{analysisbox}
            \item 当$x\to a$时，取$x_n \to a$（$x_n \neq a$），若$\lim\limits_{x\to a} f(x) = A$，则$\lim\limits_{n\to\infty} f(x_n) = A$
        \end{itemize}

        \item \textbf{应用}：
        \begin{itemize}
            \item 将数列极限转化为函数极限，利用洛必达法则等工具求解
            \item 证明函数极限不存在：找两个趋于$x_0$的数列$\{x_n\}$和$\{y_n\}$，使得$\lim\limits_{n\to\infty} f(x_n) \neq \lim\limits_{n\to\infty} f(y_n)$
        \end{itemize}
    \end{enumerate}

    \subsection{夹逼准则}

    \begin{enumerate}
        \item \textbf{定理}：若存在正整数$N$，使得当$n > N$时，有
        \[
            y_n \le x_n \le z_n
        \]
        且$\lim\limits_{n\to\infty} y_n = \lim\limits_{n\to\infty} z_n = a$，则$\lim\limits_{n\to\infty} x_n = a$

        \item \textbf{适用情形}：
        \begin{itemize}
            \item 求和或求积形式的极限
            \item 无法直接求出极限，但可估计上下界的情形
            \item 含$n$项和或$n$项积的表达式
        \end{itemize}

        \item \textbf{关键技巧}：找到合适的放缩，使两端极限相等且易求
    \end{enumerate}

    \subsubsection{放缩常用方法}
    \begin{enumerate}
        \item \textbf{利用最大最小项放缩}：
        \begin{itemize}
            \item $n\,u_{\min}\le \sum\limits_{k=1}^n u_k \le n\,u_{\max}$
            \item 若 $0\le u_k\le M\ (k=1,\dots,n)$，则$0\le \sum\limits_{k=1}^n u_k\le nM$
        \end{itemize}
        \item 利用重要不等式
        \begin{itemize}
            \item 设$a,b$为实数，则
            \begin{circlenum}
                \item 三角不等式：$|a\pm b|\le |a|+|b|$
                \item 反三角不等式：$\bigl||a|-|b|\bigr|\le |a-b|$
                \item 推广：$|a_1+a_2+\cdots+a_n| \le |a_1|+|a_2|+\cdots+|a_n|$
            \end{circlenum}
            \item 基本不等式（$a,b\geq0$）：
            \[
                \sqrt{ab}\le \frac{a+b}{2}\le \sqrt{\frac{a^2+b^2}{2}}.
            \]
            \item $|ab| \leq \displaystyle\frac{a^2 + b^2}{2}$
            \item 三元均值不等式（$a,b,c \geq 0$）：
            \[
                \sqrt[3]{abc}\leq \frac{a+b+c}{3} \leq \sqrt{\frac{a^2+b^2+c^2}{3}}.
            \]
            \item 设 $a\ge b\ge0$，则
            \begin{circlenum}
                \item 当 $m>0$ 时，$a^m\ge b^m$
                \item 当 $m<0$ 时，$a^m\le b^m$
            \end{circlenum}
            \item 若 $0<a<x<b,\ 0<c<y<d$，则$\dfrac{a}{d}<\dfrac{x}{y}<\dfrac{b}{c}$
            \item 常用初等不等式：
            \begin{circlenum}
                \item $\dfrac{1}{n^2+n} \le \dfrac{1}{n^2+k} \le \dfrac{1}{n^2+1}$（$k = 1, 2, \cdots, n$）
                \item $n \le \sqrt{n^2+k} \le \sqrt{n^2+n}$（$k = 1, 2, \cdots, n$）
                \item $0<x<\dfrac{\pi}{2}$ 时，$\dfrac{2}{\pi}x<\sin x<x<\tan x$
                \item $x > 0$时，$\sin x < x$。\emph{e.g.} $x_n>0,\ x_{n+1}=\sin x_n< x_n \Rightarrow \{x_n\}\text{单调减少}$
                \item $0 < x < \dfrac{\pi}{4}$时，$x < \tan x < \dfrac{4}{\pi}x$
                \item $x>0$ 时，$\arctan x\le x<\arcsin x$。\emph{e.g.} $x_n>0,\ x_{n+1}=\arctan x_n< x_n \Rightarrow \{x_n\}\text{单调减少}$
                \item $e^x\ge 1+x\ (x\in\mathbb{R})$
                \item $\ln x\le x-1\ (x>0)$。\emph{e.g.} 当$x_n>0$时，若$x_{n+1}=\ln x_n+1\le x_n.$，得$x_{n+1} \leq \x_n$，即$\{x_n\}$单调不增
                \item $x>0$ 时，
                \[
                    \frac{1}{1+x}<\ln\!\left(1+\frac1x\right)<\frac1x,
                    \quad\text{或}\quad
                    \frac{x}{1+x}<\ln(1+x)<x \quad \\
                \]
                \begin{analysisbox}
                    \[
                        \ln\!\left(1+\frac1x\right) = \ln\!\left(1+\frac1x\right) - \ln 1 = \ln (1+t)\big|_0^{\frac{1}{x}} = \int_0^{\frac{1}{x}} \frac{1}{1+t}\,\mathrm{d}t
                    \]
                \end{analysisbox}
            \end{circlenum}
            \item \textbf{重要极限的放缩}：若$a_1, a_2, \cdots, a_m \ge 0$，则
            \[
                \lim_{n\to\infty} \sqrt[n]{a_1^n + a_2^n + \cdots + a_m^n} = \max\{a_1, a_2, \cdots, a_m\}
            \]
        \end{itemize}
        \item 利用闭区间上连续函数必有最大值与最小值
        \item 利用压缩映射原理: 使用时，需写出证明过程
        \begin{itemize}

            \item \textbf{原理一（压缩不等式）}
            设数列 $\{x_n\}$，若存在常数 $k\,({\color{red}{0 < k < 1}})$，使得对一切
            $n=1,2,\dots$，
            \[
                |x_{n+1}-a|\le k|x_n-a|,
            \]
            则数列 $\{x_n\}$ 收敛，且
            \[
                \lim_{n\to\infty}x_n=a.
            \]

            \begin{analysisbox}[证]
                由已知条件，
                \[
                    0\le |x_{n+1}-a|\le k|x_n-a|.
                \]
                递推可得
                \[
                    \begin{aligned}
                        0\le |x_{n+1}-a|
                        &\le k|x_n-a| \\
                        &\le k^2|x_{n-1}-a| \\
                        &\ \ \vdots \\
                        &\le k^n|x_1-a|.
                    \end{aligned}
                \]
                又因为 $0<k<1$，
                \[
                    \lim_{n\to\infty}k^n=0.
                \]
                由夹逼准则，
                \[
                    \lim_{n\to\infty}|x_{n+1}-a|=0,
                \]
                从而
                \[
                    \lim_{n\to\infty}x_n=a.
                \]
            \end{analysisbox}

            \item \textbf{原理二（迭代函数收敛原理）}
            设数列 $\{x_n\}$ 由
            \[
                x_{n+1}=f(x_n), \quad n=1,2,\dots
            \]
            定义，其中 $f(x)$ 在 $\mathbb{R}$ 上可导，且方程
            \[
                f(x)=x
            \]
            有唯一解 $a$。
            若存在常数 $k\,(0<k<1)$，使得对任意 $x\in\mathbb{R}$，
            \[
                |f'(x)|\le k,
            \]
            则数列 $\{x_n\}$ 收敛，且
            \[
                \lim_{n\to\infty}x_n=a.
            \]

            \begin{analysisbox}[证]
                由于 $f(a)=a$，有
                \[
                    |x_{n+1}-a|=|f(x_n)-f(a)|.
                \]
                由拉格朗日中值定理，存在 $\xi$ 介于 $x_n$ 与 $a$ 之间，使得
                \[
                    |f(x_n)-f(a)|
                    =|f'(\xi)|\,|x_n-a|.
                \]
                又由条件 $|f'(x)|\le k<1$，可得
                \[
                    |x_{n+1}-a|\le k|x_n-a|.
                \]
                由原理一（压缩不等式），得
                \[
                    \lim_{n\to\infty}x_n=a.
                \]
            \end{analysisbox}
        \end{itemize}
    \end{enumerate}

    \subsection{单调有界准则}

    \begin{enumerate}
        \item \textbf{定理}：若数列$\{x_n\}$单调增加（减少）且有上界（下界），则$\lim\limits_{n\to\infty} x_n$存在
        \begin{itemize}
            \item 单调递增且有上界 $\Rightarrow$ 极限存在
            \item 单调递减且有下界 $\Rightarrow$ 极限存在
        \end{itemize}

        \item \textbf{应用步骤}：
        \begin{enumerate}
            \item 证明数列单调（递增或递减）
            \item 证明数列有界（上界或下界）
            \item 设$\lim\limits_{n\to\infty} x_n = A$，在递推式两边取极限求$A$
        \end{enumerate}

        \item \textbf{典型题型}：给定递推公式$x_{n+1} = f(x_n)$，证明极限存在并求极限值
    \end{enumerate}

    \subsubsection{证明数列$\{x_n\}$单调性的常用方法}
    \begin{enumerate}
        \item \textbf{差分法}：计算$x_{n+1} - x_n$
        \begin{itemize}
            \item 若$x_{n+1} - x_n > 0$，则单调递增
            \item 若$x_{n+1} - x_n < 0$，则单调递减
        \end{itemize}

        \item \textbf{比值法}：计算$\dfrac{x_{n+1}}{x_n}$（要求$x_n > 0$）
        \begin{itemize}
            \item 若$\dfrac{x_{n+1}}{x_n} > 1$，则单调递增
            \item 若$\dfrac{x_{n+1}}{x_n} < 1$，则单调递减
        \end{itemize}

        \item \textbf{数学归纳法}：
        \begin{enumerate}
            \item 验证$n=1$时不等式成立
            \item 假设$n=k$时成立
            \item 证明$n=k+1$时也成立
        \end{enumerate}
        \item \textbf{利用重要不等式}
        \item \textbf{利用函数单调性}（递推形式$x_{n+1} = f(x_n)$）：

        设$x_n \in I$，$f(x)$在区间$I$上可导：
        \begin{itemize}
            \item 若$f'(x) > 0$，则$\{x_n\}$单调
            \begin{circlenum}
                \item $x_2 > x_1 \Rightarrow \{x_n\}$单调递增
                \item $x_2 < x_1 \Rightarrow \{x_n\}$单调递减
            \end{circlenum}
            \item 若$f'(x) < 0$，则$\{x_n\}$\textbf{不单调}
        \end{itemize}
    \end{enumerate}

    \subsection{${x_n}$收敛于$a$的速度问题}

    \begin{enumerate}
        \item 设数列$\{x_n\}, \{y_n\}$在$n\to \infty$的过程中同时趋于$a$，记$u_n = |x_n - a|, v_n = |y_n - a|$，且当$n\to \infty$时，$u_n$和$v_n$都是无穷小量。若$I = \lim\limits_{n\to\infty} \displaystyle\frac{u_n}{v_n}$存在，则有

        \begin{itemize}
            \item \textbf{高阶无穷小}：当$I = 0$时，称$u_n$是$v_n$的高阶无穷小，记作$u_n = o(v_n)$。等价地，$\{x_n\}$比$\{y_n\}$更快收敛于$a$
            \item \textbf{同阶无穷小}：当$I = c \neq 0$时（$c$为常数），称$u_n$与$v_n$同阶。等价地，$\{x_n\}$与$\{y_n\}$以相同阶速度收敛于$a$
            \item \textbf{等价无穷小}：当$I = 1$时，称$u_n$与$v_n$是等价无穷小，记作$u_n \sim v_n$。这是同阶无穷小的特例
            \item \textbf{低阶无穷小}：当$I = +\infty$时，称$u_n$是$v_n$的低阶无穷小（或$v_n$是$u_n$的高阶无穷小）。等价地，$\{x_n\}$比$\{y_n\}$更慢收敛于$a$
        \end{itemize}

        \begin{analysisbox}[示例]
            设$x_n = \dfrac{1}{n}$，$y_n = \dfrac{1}{\sqrt{n}}$，$\lim\limits_{n\to \infty} x_n = \lim\limits_{n\to \infty} y_n = 0$，则
            \[
                I = \lim\limits_{n\to\infty} \dfrac{\left|\dfrac{1}{n} - 0\right|}{\left|\dfrac{1}{\sqrt{n}} - 0\right|}
                = \lim\limits_{n\to\infty} \dfrac{\dfrac{1}{n}}{\dfrac{1}{\sqrt{n}}}
                = \lim\limits_{n\to\infty} \dfrac{1}{\sqrt{n}}
                = 0
            \]
            故$\dfrac{1}{n} = o\!\left(\dfrac{1}{\sqrt{n}}\right)$，即$x_n = \dfrac{1}{n}$比$y_n = \dfrac{1}{\sqrt{n}}$更快收敛于$0$。
        \end{analysisbox}

        \item \textbf{常见收敛速度比较}（$n\to\infty$）：
        一般地更快趋于$0$：$\dfrac{1}{2^n} \gg \dfrac{1}{n!} \gg \dfrac{1}{n^p} \gg \dfrac{1}{n} \gg \dfrac{1}{\ln n}$（$p > 1$）
        \item $\color{red}{\bigstar}\color{red}{\bigstar}$数列极限定义中不体现收敛速度
        \begin{itemize}
            \item \textbf{本质区分}：极限定义是
            \[
                \forall \epsilon>0,\ \exists N,\ n>N\Rightarrow |x_n-a|<\epsilon,
            \]
            它只说明“最终进入任意邻域”，并未规定“进入得多快”。
            \item \textbf{速度信息在哪里}：速度对应的是 $N(\epsilon)$ 随 $\epsilon\to0^+$ 的增长规模；这属于定义之外的比较问题。
            \item \textbf{防混淆例子}：
            \[
                x_n=\frac1n,\qquad y_n=\frac1{\ln n}\ (n\ge2),
            \]
            二者都满足 $\lim\limits_{n\to\infty}=0$，但
            \[
                \frac1n<\epsilon \iff n>\frac1\epsilon,\qquad
                \frac1{\ln n}<\epsilon \iff n>e^{1/\epsilon}.
            \]
            同一 $\epsilon$ 下，后者所需阈值远更大，故收敛更慢。
            \item 极限定义中$\epsilon$的含义：$\epsilon$必须是不依赖于$n$的任意小正数。若$\epsilon$与$n$有关，相当于对收敛速度提出了额外要求，而极限定义并不限制收敛速度。下面的示例说明了在推理过程中，$\epsilon$ 必须作为一个与 $n$ 无关的常数来使用。
            \begin{analysisbox}[示例]
                \[
                    \begin{aligned}
                        \lim_{n \to \infty} a_n = a
                        &\Rightarrow \forall \epsilon > 0,\ \exists N > 0,\ n > N:\ |a_n - a| < \epsilon \\[4pt]
                        &\Rightarrow \bigl||a_n| - |a|\bigr| \le |a_n - a| < \epsilon
                        \quad (\text{由绝对值不等式}) \\[4pt]
                        &\Rightarrow \lim_{n \to \infty} |a_n| = |a|. \\[4pt]
                        &\Rightarrow |a| - \epsilon < |a_n| < |a| + \epsilon \\[4pt]
                    \end{aligned}
                \]
            \end{analysisbox}
            \item \textbf{深入理解：数列极限定义不体现收敛速度}
            \begin{circlenum}
                \item 极限定义只回答“是否最终进入并停留在任意小邻域内”，不回答“多快进入”。
                \item 收敛速度体现在“给定同一个$\epsilon$时，所需$N(\epsilon)$增长有多快”，这是定义之外的比较问题。
                \item 因此，两个数列都满足$\lim\limits_{n\to\infty}x_n=0$，并不意味着它们收敛快慢相同。
            \end{circlenum}
            \begin{analysisbox}[示例]
                比较
                \[
                    x_n=\frac1n,\qquad y_n=\frac1{\ln n}\ (n\ge2).
                \]
                二者都满足
                \[
                    \lim_{n\to\infty}x_n=\lim_{n\to\infty}y_n=0.
                \]
                但对同一$\epsilon\in(0,1)$，
                \[
                    \frac1n<\epsilon \iff n>\frac1\epsilon,
                    \qquad
                    \frac1{\ln n}<\epsilon \iff n>e^{1/\epsilon}.
                \]
                由于 $e^{1/\epsilon}\gg 1/\epsilon$（当$\epsilon\to0^+$），可见 $y_n$ 需要更大的 $N$ 才进入同一邻域，故 $y_n$ 收敛更慢。

                结论：极限定义本身不区分速度；速度比较需额外工具（如$o,O,\sim$、同阶/高阶无穷小）。
            \end{analysisbox}

        \end{itemize}
    \end{enumerate}

    \subsection{狄利克雷函数}

    \begin{enumerate}
        \item \textbf{定义}：
        \[
            D(x)=
            \begin{cases}
                1, & x\in\mathbb{Q},\\
                0, & x\notin\mathbb{Q}.
            \end{cases}
        \]
        \item \textbf{基本性质}：在任意区间内都同时含有有理数与无理数，因此 $D(x)$ 在任意邻域内都在 $0,1$ 之间跳动。
        \item \textbf{极限与连续性结论}：对任意 $x_0\in\mathbb{R}$，$\lim\limits_{x\to x_0}D(x)$不存在，故 $D(x)$ 在任意点都不连续。
        \item \textbf{常用判定（海涅定理）}：取
        \[
            x_n=x_0+\frac{\sqrt{2}}{n}\ (\text{无理数列}),\qquad
            y_n=x_0+\frac{1}{n}\ (\text{有理数列}),
        \]
        则 $x_n\to x_0,\ y_n\to x_0$，但
        \[
            D(x_n)=0,\qquad D(y_n)=1.
        \]
        两列函数值极限不同，故 $\lim\limits_{x\to x_0}D(x)$不存在。
    \end{enumerate}

    \subsection{结论}

    \begin{enumerate}
        \item $\lim\limits_{n\to\infty} a_n = A \; {\color[rgb]{0.2, 0.6, 0.3}{\Rightarrow}} \lim\limits_{n\to\infty} |a_n| = |A|$
        \item $\lim\limits_{x\to x_0} f(x) = A \; {\color[rgb]{0.2, 0.6, 0.3}{\Rightarrow}} \lim\limits_{x\to x_0} |f(x)| = |A|$
        \item $\lim\limits_{x\to x_0} f(x) = 0 \; {\color{red}{\Leftrightarrow}} \lim\limits_{x\to x_0} |f(x)| = 0$
        \item $\lim\limits_{n\to\infty} a_n = 0 \; {\color{red}{\Leftrightarrow}} \lim\limits_{n\to\infty} |a_n| = 0$
        \begin{analysisbox}[Example]
            若要证明$\lim\limits_{n\to\infty} a_n = 0$，只需证明
            \[
                \lim_{n\to\infty} |a_n| = 0.
            \]

            又由于 $|a_n|\ge 0$，
            可利用夹逼准则：
            若存在数列 $\{b_n\}$，使得
            \[
                0 \le |a_n| \le b_n,
                \qquad
                \lim_{n\to\infty} b_n = 0,
            \]
            则有
            \[
                \lim_{n\to\infty} a_n = 0.
            \]
        \end{analysisbox}
        \item $\color{red}{\bigstar}$若 $a_1,a_2,\dots,a_m \ge 0$，则
        \[
            \lim_{n\to\infty} \sqrt[n]{a_1^n + a_2^n + \cdots + a_m^n}
            = \max\{a_1,a_2,\dots,a_m\}
        \]
    \end{enumerate}

    \subsection{方法步骤}

    \begin{enumerate}
        \item 判断数列发散的方法
        \begin{itemize}
            \item 对于一个数列$\{a_n\}$，如果能找到一个发散的子列，则原数列一定发散
            \item 对于一个数列$\{a_n\}$，如果能找到至少两个收敛的子列$\{x_{n_k}\}$和$\{x_{n'_k}\}$，但它们收敛到不同极限，则原数列一定发散
        \end{itemize}
    \end{enumerate}

    \subsection{条件转换思路}

    \begin{enumerate}
        \item 数列$\{a_n\}$非负有界 $\Rightarrow$ 存在$0 \leq m < M$，使得对任意$n \in \mathbb{N}^+$，有$m \leq a_n \leq M$
    \end{enumerate}

\end{document}
